% Options for packages loaded elsewhere
\PassOptionsToPackage{unicode}{hyperref}
\PassOptionsToPackage{hyphens}{url}
%
\documentclass[
  11pt,
  a4paper,
]{article}
\providecommand{\tightlist}{\setlength{\itemsep}{0pt}\setlength{\parskip}{0pt}}
\usepackage{amsthm}
\usepackage{amsfonts}
\usepackage{mathrsfs}
\usepackage{mathtools}
\usepackage{stmaryrd}
\usepackage{bbm}
\usepackage{enumitem}
\usepackage{slashed}
\usepackage{tikz-cd}
\usepackage{tikz}
\usepackage{cancel}
\providecommand{\qedsymbol}{\ensuremath{\square}}
\providecommand{\real}{\mathbb{R}}
\providecommand{\Real}{\mathbb{R}}
\providecommand{\Tr}{\mathrm{Tr}}
\providecommand{\Sym}{\mathrm{Sym}}
\providecommand{\Spec}{\mathrm{Spec}}
\providecommand{\Hom}{\mathrm{Hom}}
\providecommand{\End}{\mathrm{End}}
\providecommand{\Aut}{\mathrm{Aut}}
\providecommand{\Gal}{\mathrm{Gal}}
\providecommand{\Mat}{\mathrm{Mat}}
\providecommand{\Pic}{\mathrm{Pic}}
\providecommand{\NS}{\mathrm{NS}}
\providecommand{\rank}{\mathrm{rank}}
\providecommand{\disc}{\mathrm{disc}}
\providecommand{\sgn}{\mathrm{sgn}}
\providecommand{\Lie}{\mathrm{Lie}}
\usepackage{amsmath,amssymb}
\usepackage{iftex}
\ifPDFTeX
  \usepackage[T1]{fontenc}
  \usepackage[utf8]{inputenc}
  \usepackage{textcomp} % provide euro and other symbols
\else % if luatex or xetex
  \usepackage{unicode-math} % this also loads fontspec
  \defaultfontfeatures{Scale=MatchLowercase}
  \defaultfontfeatures[\rmfamily]{Ligatures=TeX,Scale=1}
\fi
\usepackage{lmodern}
\ifPDFTeX\else
  % xetex/luatex font selection
\fi
% Use upquote if available, for straight quotes in verbatim environments
\IfFileExists{upquote.sty}{\usepackage{upquote}}{}
\IfFileExists{microtype.sty}{% use microtype if available
  \usepackage[]{microtype}
  \UseMicrotypeSet[protrusion]{basicmath} % disable protrusion for tt fonts
}{}
\makeatletter
\@ifundefined{KOMAClassName}{% if non-KOMA class
  \IfFileExists{parskip.sty}{%
    \usepackage{parskip}
  }{% else
    \setlength{\parindent}{0pt}
    \setlength{\parskip}{6pt plus 2pt minus 1pt}}
}{% if KOMA class
  \KOMAoptions{parskip=half}}
\makeatother
\usepackage{xcolor}
\usepackage[margin=1in]{geometry}
\usepackage{longtable,booktabs,array}
\usepackage{calc} % for calculating minipage widths
% Correct order of tables after \paragraph or \subparagraph
\usepackage{etoolbox}
\makeatletter
\patchcmd\longtable{\par}{\if@noskipsec\mbox{}\fi\par}{}{}
\makeatother
% Allow footnotes in longtable head/foot
\IfFileExists{footnotehyper.sty}{\usepackage{footnotehyper}}{\usepackage{footnote}}
\makesavenoteenv{longtable}
\setlength{\emergencystretch}{3em} % prevent overfull lines
\providecommand{\tightlist}{%
  \setlength{\itemsep}{0pt}\setlength{\parskip}{0pt}}
\setcounter{secnumdepth}{-\maxdimen} % remove section numbering
\ifLuaTeX
  \usepackage{selnolig}  % disable illegal ligatures
\fi
\IfFileExists{bookmark.sty}{\usepackage{bookmark}}{\usepackage{hyperref}}
\IfFileExists{xurl.sty}{\usepackage{xurl}}{} % add URL line breaks if available
\urlstyle{same}
\hypersetup{
  hidelinks,
  pdfauthor={Kévin Remondière},
  pdftitle={Formalization of the Mumford-Tate Torus Framework: Algebraicity of Hodge Cycles on Self-Products of h_K=1 CM Elliptic Curves},
  pdfcreator={LaTeX via pandoc}}

\title{Formalization of the Mumford--Tate Torus Framework: Algebraicity of Hodge Cycles on Self-Products of $h_K = 1$ CM Elliptic Curves}
\author{Kévin Remondière\thanks{Independent Researcher, Tarbes, France. ORCID: \href{https://orcid.org/0009-0008-2443-7166}{0009-0008-2443-7166}. Email: \texttt{kevin.remondiere@gmail.com}.}}
\date{May 11, 2026}

\DeclareUnicodeCharacter{2248}{\ensuremath{\approx}}
\DeclareUnicodeCharacter{2243}{\ensuremath{\simeq}}
\DeclareUnicodeCharacter{2245}{\ensuremath{\cong}}
\DeclareUnicodeCharacter{2260}{\ensuremath{\neq}}
\DeclareUnicodeCharacter{2264}{\ensuremath{\leq}}
\DeclareUnicodeCharacter{2265}{\ensuremath{\geq}}
\DeclareUnicodeCharacter{2261}{\ensuremath{\equiv}}
\DeclareUnicodeCharacter{2282}{\ensuremath{\subset}}
\DeclareUnicodeCharacter{2283}{\ensuremath{\supset}}
\DeclareUnicodeCharacter{2286}{\ensuremath{\subseteq}}
\DeclareUnicodeCharacter{2287}{\ensuremath{\supseteq}}
\DeclareUnicodeCharacter{2208}{\ensuremath{\in}}
\DeclareUnicodeCharacter{2209}{\ensuremath{\notin}}
\DeclareUnicodeCharacter{220B}{\ensuremath{\ni}}
\DeclareUnicodeCharacter{222A}{\ensuremath{\cup}}
\DeclareUnicodeCharacter{2229}{\ensuremath{\cap}}
\DeclareUnicodeCharacter{2205}{\ensuremath{\emptyset}}
\DeclareUnicodeCharacter{2200}{\ensuremath{\forall}}
\DeclareUnicodeCharacter{2203}{\ensuremath{\exists}}
\DeclareUnicodeCharacter{2202}{\ensuremath{\partial}}
\DeclareUnicodeCharacter{2207}{\ensuremath{\nabla}}
\DeclareUnicodeCharacter{221E}{\ensuremath{\infty}}
\DeclareUnicodeCharacter{2227}{\ensuremath{\wedge}}
\DeclareUnicodeCharacter{2228}{\ensuremath{\vee}}
\DeclareUnicodeCharacter{00AC}{\ensuremath{\neg}}
\DeclareUnicodeCharacter{21D2}{\ensuremath{\Rightarrow}}
\DeclareUnicodeCharacter{21D4}{\ensuremath{\Leftrightarrow}}
\DeclareUnicodeCharacter{2192}{\ensuremath{\to}}
\DeclareUnicodeCharacter{2190}{\ensuremath{\leftarrow}}
\DeclareUnicodeCharacter{2194}{\ensuremath{\leftrightarrow}}
\DeclareUnicodeCharacter{21A6}{\ensuremath{\mapsto}}
\DeclareUnicodeCharacter{2295}{\ensuremath{\oplus}}
\DeclareUnicodeCharacter{2297}{\ensuremath{\otimes}}
\DeclareUnicodeCharacter{22C5}{\ensuremath{\cdot}}
\DeclareUnicodeCharacter{00D7}{\ensuremath{\times}}
\DeclareUnicodeCharacter{00F7}{\ensuremath{\div}}
\DeclareUnicodeCharacter{00B1}{\ensuremath{\pm}}
\DeclareUnicodeCharacter{2213}{\ensuremath{\mp}}
\DeclareUnicodeCharacter{221A}{\ensuremath{\surd}}
\DeclareUnicodeCharacter{2211}{\ensuremath{\sum}}
\DeclareUnicodeCharacter{220F}{\ensuremath{\prod}}
\DeclareUnicodeCharacter{222B}{\ensuremath{\int}}
\DeclareUnicodeCharacter{222E}{\ensuremath{\oint}}
\DeclareUnicodeCharacter{03B1}{\ensuremath{\alpha}}
\DeclareUnicodeCharacter{03B2}{\ensuremath{\beta}}
\DeclareUnicodeCharacter{03B3}{\ensuremath{\gamma}}
\DeclareUnicodeCharacter{03B4}{\ensuremath{\delta}}
\DeclareUnicodeCharacter{03B5}{\ensuremath{\varepsilon}}
\DeclareUnicodeCharacter{03F5}{\ensuremath{\epsilon}}
\DeclareUnicodeCharacter{03B6}{\ensuremath{\zeta}}
\DeclareUnicodeCharacter{03B7}{\ensuremath{\eta}}
\DeclareUnicodeCharacter{03B8}{\ensuremath{\theta}}
\DeclareUnicodeCharacter{03B9}{\ensuremath{\iota}}
\DeclareUnicodeCharacter{03BA}{\ensuremath{\kappa}}
\DeclareUnicodeCharacter{03BB}{\ensuremath{\lambda}}
\DeclareUnicodeCharacter{03BC}{\ensuremath{\mu}}
\DeclareUnicodeCharacter{03BD}{\ensuremath{\nu}}
\DeclareUnicodeCharacter{03BE}{\ensuremath{\xi}}
\DeclareUnicodeCharacter{03C0}{\ensuremath{\pi}}
\DeclareUnicodeCharacter{03C1}{\ensuremath{\rho}}
\DeclareUnicodeCharacter{03C3}{\ensuremath{\sigma}}
\DeclareUnicodeCharacter{03C4}{\ensuremath{\tau}}
\DeclareUnicodeCharacter{03C5}{\ensuremath{\upsilon}}
\DeclareUnicodeCharacter{03C6}{\ensuremath{\varphi}}
\DeclareUnicodeCharacter{03D5}{\ensuremath{\phi}}
\DeclareUnicodeCharacter{03C7}{\ensuremath{\chi}}
\DeclareUnicodeCharacter{03C8}{\ensuremath{\psi}}
\DeclareUnicodeCharacter{03C9}{\ensuremath{\omega}}
\DeclareUnicodeCharacter{0393}{\ensuremath{\Gamma}}
\DeclareUnicodeCharacter{0394}{\ensuremath{\Delta}}
\DeclareUnicodeCharacter{0398}{\ensuremath{\Theta}}
\DeclareUnicodeCharacter{039B}{\ensuremath{\Lambda}}
\DeclareUnicodeCharacter{039E}{\ensuremath{\Xi}}
\DeclareUnicodeCharacter{03A0}{\ensuremath{\Pi}}
\DeclareUnicodeCharacter{03A3}{\ensuremath{\Sigma}}
\DeclareUnicodeCharacter{03A6}{\ensuremath{\Phi}}
\DeclareUnicodeCharacter{03A8}{\ensuremath{\Psi}}
\DeclareUnicodeCharacter{03A9}{\ensuremath{\Omega}}
\DeclareUnicodeCharacter{211D}{\ensuremath{\mathbb{R}}}
\DeclareUnicodeCharacter{2124}{\ensuremath{\mathbb{Z}}}
\DeclareUnicodeCharacter{211A}{\ensuremath{\mathbb{Q}}}
\DeclareUnicodeCharacter{2102}{\ensuremath{\mathbb{C}}}
\DeclareUnicodeCharacter{2115}{\ensuremath{\mathbb{N}}}
\DeclareUnicodeCharacter{2119}{\ensuremath{\mathbb{P}}}
\DeclareUnicodeCharacter{210B}{\ensuremath{\mathcal{H}}}
\DeclareUnicodeCharacter{2112}{\ensuremath{\mathcal{L}}}
\DeclareUnicodeCharacter{2113}{\ensuremath{\ell}}
\DeclareUnicodeCharacter{2200}{\ensuremath{\forall}}
\DeclareUnicodeCharacter{1D4AA}{\ensuremath{\mathcal{O}}}
\DeclareUnicodeCharacter{1D52D}{\ensuremath{\mathfrak{p}}}
\DeclareUnicodeCharacter{1D52A}{\ensuremath{\mathfrak{m}}}
\DeclareUnicodeCharacter{1D52B}{\ensuremath{\mathfrak{n}}}
\DeclareUnicodeCharacter{1D51E}{\ensuremath{\mathfrak{a}}}
\DeclareUnicodeCharacter{1D51F}{\ensuremath{\mathfrak{b}}}
\DeclareUnicodeCharacter{1D52E}{\ensuremath{\mathfrak{q}}}
\DeclareUnicodeCharacter{1D516}{\ensuremath{\mathfrak{S}}}
\DeclareUnicodeCharacter{1D53D}{\ensuremath{\mathbb{F}}}
\DeclareUnicodeCharacter{00B0}{\ensuremath{^\circ}}
\DeclareUnicodeCharacter{2032}{\ensuremath{{}'}}
\DeclareUnicodeCharacter{2033}{\ensuremath{{}''}}
\DeclareUnicodeCharacter{00B7}{\ensuremath{\cdot}}
\DeclareUnicodeCharacter{2026}{\ensuremath{\ldots}}
\DeclareUnicodeCharacter{2010}{-}
\DeclareUnicodeCharacter{2013}{--}
\DeclareUnicodeCharacter{2014}{---}
\DeclareUnicodeCharacter{2018}{`}
\DeclareUnicodeCharacter{2019}{'}
\DeclareUnicodeCharacter{201C}{``}
\DeclareUnicodeCharacter{201D}{''}
\DeclareUnicodeCharacter{00A0}{~}
\DeclareUnicodeCharacter{2009}{\,}
\DeclareUnicodeCharacter{200A}{\,}
\DeclareUnicodeCharacter{2002}{\,}
\DeclareUnicodeCharacter{2003}{\;}
\DeclareUnicodeCharacter{2020}{\dag}
\DeclareUnicodeCharacter{2021}{\ddag}
\DeclareUnicodeCharacter{2308}{\ensuremath{\lceil}}
\DeclareUnicodeCharacter{2309}{\ensuremath{\rceil}}
\DeclareUnicodeCharacter{230A}{\ensuremath{\lfloor}}
\DeclareUnicodeCharacter{230B}{\ensuremath{\rfloor}}
\DeclareUnicodeCharacter{27E8}{\ensuremath{\langle}}
\DeclareUnicodeCharacter{27E9}{\ensuremath{\rangle}}
\DeclareUnicodeCharacter{2225}{\ensuremath{\parallel}}
\DeclareUnicodeCharacter{2226}{\ensuremath{\nparallel}}
\DeclareUnicodeCharacter{2218}{\ensuremath{\circ}}
\DeclareUnicodeCharacter{2236}{\ensuremath{\colon}}
\DeclareUnicodeCharacter{2237}{\ensuremath{::}}
\DeclareUnicodeCharacter{2254}{\ensuremath{:=}}
\DeclareUnicodeCharacter{2255}{\ensuremath{=:}}
\DeclareUnicodeCharacter{2A7D}{\ensuremath{\leqslant}}
\DeclareUnicodeCharacter{2A7E}{\ensuremath{\geqslant}}
\DeclareUnicodeCharacter{226A}{\ensuremath{\ll}}
\DeclareUnicodeCharacter{226B}{\ensuremath{\gg}}
\DeclareUnicodeCharacter{2272}{\ensuremath{\lesssim}}
\DeclareUnicodeCharacter{2273}{\ensuremath{\gtrsim}}
\DeclareUnicodeCharacter{2329}{\ensuremath{\langle}}
\DeclareUnicodeCharacter{232A}{\ensuremath{\rangle}}
\DeclareUnicodeCharacter{25A1}{\ensuremath{\square}}
\DeclareUnicodeCharacter{25E6}{\ensuremath{\circ}}
\DeclareUnicodeCharacter{2660}{\ensuremath{\spadesuit}}
\DeclareUnicodeCharacter{2663}{\ensuremath{\clubsuit}}
\DeclareUnicodeCharacter{2665}{\ensuremath{\heartsuit}}
\DeclareUnicodeCharacter{2666}{\ensuremath{\diamondsuit}}
\DeclareUnicodeCharacter{2713}{\ensuremath{\checkmark}}
\DeclareUnicodeCharacter{2717}{\textsf{x}}
\DeclareUnicodeCharacter{00BD}{\ensuremath{\tfrac{1}{2}}}
\DeclareUnicodeCharacter{00BC}{\ensuremath{\tfrac{1}{4}}}
\DeclareUnicodeCharacter{00BE}{\ensuremath{\tfrac{3}{4}}}
\DeclareUnicodeCharacter{00A7}{\S}
\DeclareUnicodeCharacter{00B6}{\P}
\DeclareUnicodeCharacter{2022}{\textbullet}
\DeclareUnicodeCharacter{2023}{\ensuremath{\blacktriangleright}}
\DeclareUnicodeCharacter{20AC}{\texteuro}
\DeclareUnicodeCharacter{00A3}{\textsterling}
\DeclareUnicodeCharacter{2212}{\ensuremath{-}}
\DeclareUnicodeCharacter{220E}{\ensuremath{\blacksquare}}
\DeclareUnicodeCharacter{2070}{\ensuremath{{}^{0}}}
\DeclareUnicodeCharacter{00B9}{\ensuremath{{}^{1}}}
\DeclareUnicodeCharacter{00B2}{\ensuremath{{}^{2}}}
\DeclareUnicodeCharacter{00B3}{\ensuremath{{}^{3}}}
\DeclareUnicodeCharacter{2074}{\ensuremath{{}^{4}}}
\DeclareUnicodeCharacter{2075}{\ensuremath{{}^{5}}}
\DeclareUnicodeCharacter{2076}{\ensuremath{{}^{6}}}
\DeclareUnicodeCharacter{2077}{\ensuremath{{}^{7}}}
\DeclareUnicodeCharacter{2078}{\ensuremath{{}^{8}}}
\DeclareUnicodeCharacter{2079}{\ensuremath{{}^{9}}}
\DeclareUnicodeCharacter{2080}{\ensuremath{{}_{0}}}
\DeclareUnicodeCharacter{2081}{\ensuremath{{}_{1}}}
\DeclareUnicodeCharacter{2082}{\ensuremath{{}_{2}}}
\DeclareUnicodeCharacter{2083}{\ensuremath{{}_{3}}}
\DeclareUnicodeCharacter{2084}{\ensuremath{{}_{4}}}
\DeclareUnicodeCharacter{2085}{\ensuremath{{}_{5}}}
\DeclareUnicodeCharacter{2086}{\ensuremath{{}_{6}}}
\DeclareUnicodeCharacter{2087}{\ensuremath{{}_{7}}}
\DeclareUnicodeCharacter{2088}{\ensuremath{{}_{8}}}
\DeclareUnicodeCharacter{2089}{\ensuremath{{}_{9}}}
\DeclareUnicodeCharacter{207A}{\ensuremath{{}^{+}}}
\DeclareUnicodeCharacter{207B}{\ensuremath{{}^{-}}}
\DeclareUnicodeCharacter{207C}{\ensuremath{{}^{=}}}
\DeclareUnicodeCharacter{207D}{\ensuremath{{}^{(}}}
\DeclareUnicodeCharacter{207E}{\ensuremath{{}^{)}}}
\DeclareUnicodeCharacter{208A}{\ensuremath{{}_{+}}}
\DeclareUnicodeCharacter{208B}{\ensuremath{{}_{-}}}
\DeclareUnicodeCharacter{208C}{\ensuremath{{}_{=}}}
\DeclareUnicodeCharacter{208D}{\ensuremath{{}_{(}}}
\DeclareUnicodeCharacter{208E}{\ensuremath{{}_{)}}}
\DeclareUnicodeCharacter{0300}{\`{}}
\DeclareUnicodeCharacter{0301}{\'{}}
\DeclareUnicodeCharacter{0302}{\^{}}
\DeclareUnicodeCharacter{0303}{\~{}}
\DeclareUnicodeCharacter{0304}{\={}}
\DeclareUnicodeCharacter{0306}{\u{}}
\DeclareUnicodeCharacter{0307}{\.{}}
\DeclareUnicodeCharacter{0308}{\\\"{}}
\DeclareUnicodeCharacter{030A}{\r{}}
\DeclareUnicodeCharacter{030B}{\H{}}
\DeclareUnicodeCharacter{030C}{\v{}}
\DeclareUnicodeCharacter{2032}{\ensuremath{{}'}}
\DeclareUnicodeCharacter{2034}{\ensuremath{{}'''}}
\DeclareUnicodeCharacter{2057}{\ensuremath{{}''''}}
\DeclareUnicodeCharacter{2220}{\ensuremath{\angle}}
\DeclareUnicodeCharacter{2221}{\ensuremath{\measuredangle}}
\DeclareUnicodeCharacter{2222}{\ensuremath{\sphericalangle}}
\DeclareUnicodeCharacter{2223}{\ensuremath{\mid}}
\DeclareUnicodeCharacter{2224}{\ensuremath{\nmid}}
\DeclareUnicodeCharacter{2234}{\ensuremath{\therefore}}
\DeclareUnicodeCharacter{2235}{\ensuremath{\because}}
\DeclareUnicodeCharacter{29F8}{\ensuremath{/}}
\DeclareUnicodeCharacter{29F9}{\ensuremath{\backslash}}
\DeclareUnicodeCharacter{2A00}{\ensuremath{\bigodot}}
\DeclareUnicodeCharacter{2A01}{\ensuremath{\bigoplus}}
\DeclareUnicodeCharacter{2A02}{\ensuremath{\bigotimes}}
\DeclareUnicodeCharacter{2A03}{\ensuremath{\biguplus}}
\DeclareUnicodeCharacter{2A04}{\ensuremath{\biguplus}}
\DeclareUnicodeCharacter{2A0C}{\ensuremath{\iiiint}}
\DeclareUnicodeCharacter{223C}{\ensuremath{\sim}}
\DeclareUnicodeCharacter{223D}{\ensuremath{\backsim}}
\DeclareUnicodeCharacter{2240}{\ensuremath{\wr}}
\DeclareUnicodeCharacter{2249}{\ensuremath{\not\approx}}
\DeclareUnicodeCharacter{221D}{\ensuremath{\propto}}
\DeclareUnicodeCharacter{210D}{\ensuremath{\mathbb{H}}}
\DeclareUnicodeCharacter{210E}{\ensuremath{h}}
\DeclareUnicodeCharacter{210A}{\ensuremath{g}}
\DeclareUnicodeCharacter{2107}{\ensuremath{e}}
\DeclareUnicodeCharacter{2127}{\ensuremath{\mho}}
\DeclareUnicodeCharacter{2129}{\ensuremath{\iota}}
\DeclareUnicodeCharacter{2300}{\ensuremath{\diameter}}
\DeclareUnicodeCharacter{29B0}{\ensuremath{\emptyset}}
\DeclareUnicodeCharacter{2A0D}{\ensuremath{\fint}}
\DeclareUnicodeCharacter{2235}{\ensuremath{\because}}
\DeclareUnicodeCharacter{2135}{\ensuremath{\aleph}}
\DeclareUnicodeCharacter{2136}{\ensuremath{\beth}}
\DeclareUnicodeCharacter{2137}{\ensuremath{\gimel}}
\DeclareUnicodeCharacter{2138}{\ensuremath{\daleth}}
\DeclareUnicodeCharacter{22EE}{\ensuremath{\vdots}}
\DeclareUnicodeCharacter{22EF}{\ensuremath{\cdots}}
\DeclareUnicodeCharacter{22F0}{\ensuremath{\iddots}}
\DeclareUnicodeCharacter{22F1}{\ensuremath{\ddots}}
\DeclareUnicodeCharacter{207F}{\ensuremath{{}^{n}}}
\DeclareUnicodeCharacter{2071}{\ensuremath{{}^{i}}}
\DeclareUnicodeCharacter{2299}{\ensuremath{\odot}}
\DeclareUnicodeCharacter{229E}{\ensuremath{\boxplus}}
\DeclareUnicodeCharacter{229F}{\ensuremath{\boxminus}}
\DeclareUnicodeCharacter{22A0}{\ensuremath{\boxtimes}}
\DeclareUnicodeCharacter{22A2}{\ensuremath{\vdash}}
\DeclareUnicodeCharacter{22A3}{\ensuremath{\dashv}}
\DeclareUnicodeCharacter{22A4}{\ensuremath{\top}}
\DeclareUnicodeCharacter{22A5}{\ensuremath{\bot}}
\DeclareUnicodeCharacter{22A8}{\ensuremath{\models}}
\DeclareUnicodeCharacter{22CA}{\ensuremath{\rtimes}}
\DeclareUnicodeCharacter{22C9}{\ensuremath{\ltimes}}
\DeclareUnicodeCharacter{22CB}{\ensuremath{\leftthreetimes}}
\DeclareUnicodeCharacter{22CC}{\ensuremath{\rightthreetimes}}
\DeclareUnicodeCharacter{2A0F}{\ensuremath{\fint}}
\DeclareUnicodeCharacter{27F6}{\ensuremath{\longrightarrow}}
\DeclareUnicodeCharacter{27F5}{\ensuremath{\longleftarrow}}
\DeclareUnicodeCharacter{27F7}{\ensuremath{\longleftrightarrow}}
\DeclareUnicodeCharacter{27F8}{\ensuremath{\Longleftarrow}}
\DeclareUnicodeCharacter{27F9}{\ensuremath{\Longrightarrow}}
\DeclareUnicodeCharacter{27FA}{\ensuremath{\Longleftrightarrow}}
\DeclareUnicodeCharacter{27FC}{\ensuremath{\longmapsto}}
\DeclareUnicodeCharacter{2191}{\ensuremath{\uparrow}}
\DeclareUnicodeCharacter{2193}{\ensuremath{\downarrow}}
\DeclareUnicodeCharacter{21D1}{\ensuremath{\Uparrow}}
\DeclareUnicodeCharacter{21D3}{\ensuremath{\Downarrow}}
\DeclareUnicodeCharacter{21AA}{\ensuremath{\hookrightarrow}}
\DeclareUnicodeCharacter{21A9}{\ensuremath{\hookleftarrow}}
\DeclareUnicodeCharacter{21A0}{\ensuremath{\twoheadrightarrow}}
\DeclareUnicodeCharacter{219E}{\ensuremath{\twoheadleftarrow}}
\DeclareUnicodeCharacter{21A3}{\ensuremath{\rightarrowtail}}
\DeclareUnicodeCharacter{21A2}{\ensuremath{\leftarrowtail}}
\DeclareUnicodeCharacter{210F}{\ensuremath{\hbar}}
\DeclareUnicodeCharacter{2118}{\ensuremath{\wp}}
\DeclareUnicodeCharacter{2202}{\ensuremath{\partial}}
\DeclareUnicodeCharacter{1D538}{\ensuremath{\mathbb{A}}}
\DeclareUnicodeCharacter{1D539}{\ensuremath{\mathbb{B}}}
\DeclareUnicodeCharacter{1D53B}{\ensuremath{\mathbb{D}}}
\DeclareUnicodeCharacter{1D53C}{\ensuremath{\mathbb{E}}}
\DeclareUnicodeCharacter{1D53E}{\ensuremath{\mathbb{G}}}
\DeclareUnicodeCharacter{1D540}{\ensuremath{\mathbb{I}}}
\DeclareUnicodeCharacter{1D541}{\ensuremath{\mathbb{J}}}
\DeclareUnicodeCharacter{1D542}{\ensuremath{\mathbb{K}}}
\DeclareUnicodeCharacter{1D543}{\ensuremath{\mathbb{L}}}
\DeclareUnicodeCharacter{1D544}{\ensuremath{\mathbb{M}}}
\DeclareUnicodeCharacter{1D546}{\ensuremath{\mathbb{O}}}
\DeclareUnicodeCharacter{1D54A}{\ensuremath{\mathbb{S}}}
\DeclareUnicodeCharacter{1D54B}{\ensuremath{\mathbb{T}}}
\DeclareUnicodeCharacter{1D54C}{\ensuremath{\mathbb{U}}}
\DeclareUnicodeCharacter{1D54D}{\ensuremath{\mathbb{V}}}
\DeclareUnicodeCharacter{1D54E}{\ensuremath{\mathbb{W}}}
\DeclareUnicodeCharacter{1D54F}{\ensuremath{\mathbb{X}}}
\DeclareUnicodeCharacter{1D550}{\ensuremath{\mathbb{Y}}}
\DeclareUnicodeCharacter{1D504}{\ensuremath{\mathfrak{A}}}
\DeclareUnicodeCharacter{1D505}{\ensuremath{\mathfrak{B}}}
\DeclareUnicodeCharacter{212D}{\ensuremath{\mathfrak{C}}}
\DeclareUnicodeCharacter{1D507}{\ensuremath{\mathfrak{D}}}
\DeclareUnicodeCharacter{1D508}{\ensuremath{\mathfrak{E}}}
\DeclareUnicodeCharacter{1D509}{\ensuremath{\mathfrak{F}}}
\DeclareUnicodeCharacter{1D50A}{\ensuremath{\mathfrak{G}}}
\DeclareUnicodeCharacter{210C}{\ensuremath{\mathfrak{H}}}
\DeclareUnicodeCharacter{2111}{\ensuremath{\mathfrak{I}}}
\DeclareUnicodeCharacter{1D50D}{\ensuremath{\mathfrak{J}}}
\DeclareUnicodeCharacter{1D50E}{\ensuremath{\mathfrak{K}}}
\DeclareUnicodeCharacter{1D50F}{\ensuremath{\mathfrak{L}}}
\DeclareUnicodeCharacter{1D510}{\ensuremath{\mathfrak{M}}}
\DeclareUnicodeCharacter{1D511}{\ensuremath{\mathfrak{N}}}
\DeclareUnicodeCharacter{1D512}{\ensuremath{\mathfrak{O}}}
\DeclareUnicodeCharacter{1D513}{\ensuremath{\mathfrak{P}}}
\DeclareUnicodeCharacter{1D514}{\ensuremath{\mathfrak{Q}}}
\DeclareUnicodeCharacter{211C}{\ensuremath{\mathfrak{R}}}
\DeclareUnicodeCharacter{1D516}{\ensuremath{\mathfrak{S}}}
\DeclareUnicodeCharacter{1D517}{\ensuremath{\mathfrak{T}}}
\DeclareUnicodeCharacter{1D518}{\ensuremath{\mathfrak{U}}}
\DeclareUnicodeCharacter{1D519}{\ensuremath{\mathfrak{V}}}
\DeclareUnicodeCharacter{1D51A}{\ensuremath{\mathfrak{W}}}
\DeclareUnicodeCharacter{1D51B}{\ensuremath{\mathfrak{X}}}
\DeclareUnicodeCharacter{1D51C}{\ensuremath{\mathfrak{Y}}}
\DeclareUnicodeCharacter{2128}{\ensuremath{\mathfrak{Z}}}
\DeclareUnicodeCharacter{1D520}{\ensuremath{\mathfrak{c}}}
\DeclareUnicodeCharacter{1D521}{\ensuremath{\mathfrak{d}}}
\DeclareUnicodeCharacter{1D522}{\ensuremath{\mathfrak{e}}}
\DeclareUnicodeCharacter{1D523}{\ensuremath{\mathfrak{f}}}
\DeclareUnicodeCharacter{1D524}{\ensuremath{\mathfrak{g}}}
\DeclareUnicodeCharacter{1D525}{\ensuremath{\mathfrak{h}}}
\DeclareUnicodeCharacter{1D526}{\ensuremath{\mathfrak{i}}}
\DeclareUnicodeCharacter{1D527}{\ensuremath{\mathfrak{j}}}
\DeclareUnicodeCharacter{1D528}{\ensuremath{\mathfrak{k}}}
\DeclareUnicodeCharacter{1D529}{\ensuremath{\mathfrak{l}}}
\DeclareUnicodeCharacter{1D52C}{\ensuremath{\mathfrak{o}}}
\DeclareUnicodeCharacter{1D52F}{\ensuremath{\mathfrak{r}}}
\DeclareUnicodeCharacter{1D530}{\ensuremath{\mathfrak{s}}}
\DeclareUnicodeCharacter{1D531}{\ensuremath{\mathfrak{t}}}
\DeclareUnicodeCharacter{1D532}{\ensuremath{\mathfrak{u}}}
\DeclareUnicodeCharacter{1D533}{\ensuremath{\mathfrak{v}}}
\DeclareUnicodeCharacter{1D534}{\ensuremath{\mathfrak{w}}}
\DeclareUnicodeCharacter{1D535}{\ensuremath{\mathfrak{x}}}
\DeclareUnicodeCharacter{1D536}{\ensuremath{\mathfrak{y}}}
\DeclareUnicodeCharacter{1D537}{\ensuremath{\mathfrak{z}}}
\DeclareUnicodeCharacter{1D49C}{\ensuremath{\mathcal{A}}}
\DeclareUnicodeCharacter{212C}{\ensuremath{\mathcal{B}}}
\DeclareUnicodeCharacter{1D49E}{\ensuremath{\mathcal{C}}}
\DeclareUnicodeCharacter{1D49F}{\ensuremath{\mathcal{D}}}
\DeclareUnicodeCharacter{2130}{\ensuremath{\mathcal{E}}}
\DeclareUnicodeCharacter{2131}{\ensuremath{\mathcal{F}}}
\DeclareUnicodeCharacter{1D4A2}{\ensuremath{\mathcal{G}}}
\DeclareUnicodeCharacter{210B}{\ensuremath{\mathcal{H}}}
\DeclareUnicodeCharacter{2110}{\ensuremath{\mathcal{I}}}
\DeclareUnicodeCharacter{1D4A5}{\ensuremath{\mathcal{J}}}
\DeclareUnicodeCharacter{1D4A6}{\ensuremath{\mathcal{K}}}
\DeclareUnicodeCharacter{2112}{\ensuremath{\mathcal{L}}}
\DeclareUnicodeCharacter{2133}{\ensuremath{\mathcal{M}}}
\DeclareUnicodeCharacter{1D4A9}{\ensuremath{\mathcal{N}}}
\DeclareUnicodeCharacter{1D4AB}{\ensuremath{\mathcal{P}}}
\DeclareUnicodeCharacter{1D4AC}{\ensuremath{\mathcal{Q}}}
\DeclareUnicodeCharacter{211B}{\ensuremath{\mathcal{R}}}
\DeclareUnicodeCharacter{1D4AE}{\ensuremath{\mathcal{S}}}
\DeclareUnicodeCharacter{1D4AF}{\ensuremath{\mathcal{T}}}
\DeclareUnicodeCharacter{1D4B0}{\ensuremath{\mathcal{U}}}
\DeclareUnicodeCharacter{1D4B1}{\ensuremath{\mathcal{V}}}
\DeclareUnicodeCharacter{1D4B2}{\ensuremath{\mathcal{W}}}
\DeclareUnicodeCharacter{1D4B3}{\ensuremath{\mathcal{X}}}
\DeclareUnicodeCharacter{1D4B4}{\ensuremath{\mathcal{Y}}}
\DeclareUnicodeCharacter{1D4B5}{\ensuremath{\mathcal{Z}}}
\begin{document}
\maketitle

\begin{abstract}
We compute the Mumford--Tate group of the rational Hodge structure $H^1(E_K, \mathbb{Q})$ for the six imaginary quadratic discriminants $D \in \{-7, -11, -19, -43, -67, -163\}$ of class number one (excluding $D \in \{-3, -4, -8\}$ for unit-group reasons), where $E_K / \mathbb{Q}$ denotes the canonical CM elliptic curve. The group is the Weil-restriction torus $\mathrm{Res}_{K/\mathbb{Q}}\,\mathbb{G}_m$. Using the standard torus criterion of Deligne (1979) and Milne (Introduction to Shimura Varieties, 2017), we deduce that all Hodge cycles on the self-products $(E_K)^n$ are algebraic for every $n \geq 1$, and we give explicit cycle representatives in the case $n = 4$ relevant to weight-$5$ CM newforms. We also describe the transcendental lattice of the associated Kummer K3 surface $\widetilde{X}_D = \mathrm{Km}(E_K \times E_K)$, and explain how the framework supersedes a previously circulated speculative attribution in the literature.
\end{abstract}

\begin{center}\rule{0.5\linewidth}{0.5pt}\end{center}

\hypertarget{notation-conventions-and-the-six-discriminants}{%
\subsection{1. Notation, conventions, and the six discriminants}\label{notation-conventions-and-the-six-discriminants}}

Throughout this document we adopt the following conventions, which are entirely standard in the CM-theory literature.

\begin{itemize}
\tightlist
\item
  \(D \in \{-7, -11, -19, -43, -67, -163\}\) denotes one of the \textbf{six} imaginary quadratic fundamental discriminants with class number \(h_K = 1\) that admit a \emph{canonical} CM elliptic curve \(E_K / \mathbb{Q}\). The full Heegner list of \(h_K = 1\) discriminants is \(\{-3, -4, -7, -8, -11, -19, -43, -67, -163\}\) ; we exclude \(D \in \{-3, -4, -8\}\) because the unit groups \(\mathcal{O}_K^* = \mu_6, \mu_4, \mu_2\) for these three contribute extra automorphisms that complicate the Mumford-Tate calculation by a finite-quotient correction. The six retained discriminants all satisfy \(\mathcal{O}_K^* = \{\pm 1\}\), so their Mumford-Tate analysis is uniform.
\item
  \(K = K_D := \mathbb{Q}(\sqrt{D})\) is the imaginary quadratic field of discriminant \(D\) ; \(\mathcal{O}_K\) is its ring of integers. For each of the six \(D\), \(\mathcal{O}_K = \mathbb{Z}[\omega_D]\) with \(\omega_D = (1 + \sqrt{D})/2\) if \(D \equiv 1 \pmod 4\) (which holds for \(D = -7, -11, -19, -43, -67, -163\), all \(\equiv 1 \pmod 4\)).
\item
  \(E_K / \mathbb{Q}\) is the \textbf{canonical CM elliptic curve} with \(\mathrm{End}_{\overline{\mathbb{Q}}}(E_K) = \mathcal{O}_K\). Existence over \(\mathbb{Q}\) (rather than merely over the Hilbert class field \(H_K = K\) since \(h_K = 1\)) follows from Deuring's correspondence and the Heegner \(j\)-invariant being rational (\(j(E_K) \in \mathbb{Z}\), see e.g.~Silverman, \emph{Advanced Topics in the Arithmetic of Elliptic Curves}, GTM 151, 1994, Ch. II §6 ; or Cox, \emph{Primes of the Form} \(x^2 + ny^2\), 2nd ed., Wiley 2013, §13). Explicit Heegner \(j\)-values are \(j(E_{-7}) = -3375\), \(j(E_{-11}) = -32768\), \(j(E_{-19}) = -884736\), \(j(E_{-43}) = -884736000\), \(j(E_{-67}) = -147197952000\), \(j(E_{-163}) = -262537412640768000\) --- all integers, all matching the modular form values to 200+ digits.
\item
  \(\widetilde{X}_D\) denotes the (smooth, projective) \textbf{Kummer K3 surface} \(\mathrm{Km}(E_K \times E_K) = \widetilde{(E_K \times E_K) / \langle -1 \rangle}\) associated to the self-product of \(E_K\). Its Picard rank is \(\rho(\widetilde{X}_D) = 20\) (the maximal rank in characteristic 0 by Lefschetz, achieved precisely on singular K3 surfaces in the Shioda--Inose sense ; see Shioda--Inose 1977 \emph{Proc. Symp. Pure Math.} 29, 119-136, and Schütt 2008 \emph{Algebra and Number Theory} 2, 357-401, arXiv:0804.1558 --- verified REAL via cross-check). Its \textbf{transcendental lattice} \(T(\widetilde{X}_D)\) has rank 2 and motivic weight 2.
\item
  \(H^*(-, \mathbb{Q})\) denotes singular (Betti) cohomology with rational coefficients ; \(H^*(-, \mathbb{Q}_\ell)\) the étale cohomology ; we shall freely identify these via the comparison isomorphism (Artin) when discussing Hodge-theoretic statements.
\item
  For a smooth complex projective variety \(X\) of complex dimension \(d\) and integer \(n \in \{0, 1, \ldots, 2d\}\), we write \(H^n(X)_\mathbb{C} = \bigoplus_{p + q = n} H^{p, q}(X)\) for the \textbf{Hodge decomposition}. A class \(\xi \in H^n(X, \mathbb{Q})\) is called a \textbf{Hodge class of type} \((p, p)\) (for \(n = 2p\) even) if its image in \(H^n(X, \mathbb{C})\) lies in \(H^{p, p}(X)\).
\item
  The \textbf{Mumford-Tate group} \(\mathrm{MT}(V) \subset \mathrm{GL}(V_\mathbb{Q})\) of a polarizable rational Hodge structure \(V\) is the smallest \(\mathbb{Q}\)-algebraic subgroup such that the Hodge-structure cocharacter \(h : \mathbb{S} \to \mathrm{GL}(V_\mathbb{R})\) factors through \(\mathrm{MT}(V)_\mathbb{R}\) (Deligne 1979 \emph{Proc. Symp. Pure Math.} 33 part 2, §3 ; Milne 2017 \emph{Introduction to Shimura Varieties}, §6, available at https://www.jmilne.org/math/xnotes/svi.pdf --- REAL, lecture notes).
\item
  \(\mathrm{Res}_{K/\mathbb{Q}}\,\mathbb{G}_m\) denotes the \textbf{Weil restriction} of the multiplicative group from \(K\) to \(\mathbb{Q}\) ; this is a 2-dimensional algebraic torus over \(\mathbb{Q}\) with \(\mathbb{Q}\)-points \(K^*\) and \(\mathbb{R}\)-points \(\mathbb{C}^*\). The natural diagonal embedding \(\mathbb{G}_m \hookrightarrow \mathrm{Res}_{K/\mathbb{Q}}\,\mathbb{G}_m\) (corresponding to \(\mathbb{Q}^* \hookrightarrow K^*\)) gives a quotient torus \(\mathrm{Res}_{K/\mathbb{Q}}\,\mathbb{G}_m / \mathbb{G}_m\), which is 1-dimensional and \(\mathbb{Q}\)-rational ; over \(\overline{\mathbb{Q}}\) it splits as the kernel torus of the norm \(N_{K/\mathbb{Q}} : \mathrm{Res}_{K/\mathbb{Q}}\,\mathbb{G}_m \to \mathbb{G}_m\).
\end{itemize}

\begin{center}\rule{0.5\linewidth}{0.5pt}\end{center}

\hypertarget{the-mumford-tate-group-of-a-cm-elliptic-curve-and-its-symmetric-powers}{%
\subsection{2. The Mumford-Tate group of a CM elliptic curve and its symmetric powers}\label{the-mumford-tate-group-of-a-cm-elliptic-curve-and-its-symmetric-powers}}

This section is purely expository ; all results are classical (1966-1982).

\hypertarget{the-mumford-tate-group-of-e_k}{%
\subsubsection{\texorpdfstring{2.1 The Mumford-Tate group of \(E_K\)}{2.1 The Mumford-Tate group of E\_K}}\label{the-mumford-tate-group-of-e_k}}

\textbf{Theorem 2.1 (Mumford 1966 ; standard).} \emph{For each of the six \(D \in \{-7, -11, -19, -43, -67, -163\}\), the Mumford-Tate group of the rational Hodge structure \(H^1(E_K, \mathbb{Q})\) is}
\[
\mathrm{MT}(H^1(E_K, \mathbb{Q})) \;\cong\; \mathrm{Res}_{K/\mathbb{Q}}\,\mathbb{G}_m,
\]
\emph{the 2-dimensional Weil-restriction torus, where \(K = K_D = \mathbb{Q}(\sqrt{D})\).}

\textbf{Proof sketch.} \(H^1(E_K, \mathbb{Q})\) has rank 2 and carries the action of \(\mathrm{End}_{\overline{\mathbb{Q}}}(E_K) \otimes \mathbb{Q} = K\) by the Hodge-theoretic functoriality of \(K \subset \mathrm{End}^0(E_K)\). The Hodge decomposition \(H^1(E_K, \mathbb{C}) = H^{1, 0} \oplus H^{0, 1}\) is preserved by this \(K\)-action because \(K\) acts by \emph{holomorphic} endomorphisms. The Hodge cocharacter \(h : \mathbb{S} \to \mathrm{GL}(H^1)_\mathbb{R}\) therefore commutes with the \(K\)-action, so \(h\) factors through the centralizer of \(K\) in \(\mathrm{GL}_2\), which is the torus \(\mathrm{Res}_{K/\mathbb{Q}}\,\mathbb{G}_m\). By minimality of MT in the definition, \(\mathrm{MT}(H^1) \subseteq \mathrm{Res}_{K/\mathbb{Q}}\,\mathbb{G}_m\). The reverse inclusion follows because the image of \(h_\mathbb{R}\) generates a Zariski-dense subgroup of \(\mathrm{Res}_{K/\mathbb{Q}}\,\mathbb{G}_m(\mathbb{R}) = \mathbb{C}^*\) (the image is the unit circle if we restrict to special-orthogonal cocharacters, but the \emph{full} Mumford-Tate group is generated by all \(\mathrm{Aut}(\mathbb{C}/\mathbb{R})\)-conjugates and includes the full \(\mathbb{C}^*\)). See Mumford 1966 \emph{Math. Ann.} 181, 345-351, ``\emph{Families of abelian varieties}'' (REAL, classical, MathSciNet MR0207000) ; or Deligne--Milne--Ogus--Shih 1982 \emph{LNM} 900, Ch. I (Deligne) ``Hodge cycles on abelian varieties'', §3 --- REAL, Springer LNM 900, pp.~9-100, MathSciNet MR0654325. \(\blacksquare\)

\textbf{Remark 2.2.} The statement ``MT is a torus'' is the \emph{single} algebraic-geometric input that makes the six \(h_K = 1\) discriminants tractable in the present framework. For non-CM elliptic curves, \(\mathrm{MT}(H^1)\) is the \emph{full} \(\mathrm{GL}_2\) or \(\mathrm{SL}_2 \times \mathbb{G}_m\) depending on convention, and Hodge-cycle algebraicity becomes a deep open problem (the Mumford-Tate conjecture is itself open in general). For the CM case, by Theorem 2.1, MT is a torus, and Theorem 3.3 below gives unconditional algebraicity of all Hodge cycles.

\hypertarget{the-mumford-tate-group-of-e_kn}{%
\subsubsection{\texorpdfstring{2.2 The Mumford-Tate group of \((E_K)^n\)}{2.2 The Mumford-Tate group of (E\_K)\^{}n}}\label{the-mumford-tate-group-of-e_kn}}

\textbf{Corollary 2.3 (multiplicativity).} \emph{For each integer \(n \geq 1\),}
\[
\mathrm{MT}\bigl(H^1((E_K)^n, \mathbb{Q})\bigr) \;=\; \mathrm{MT}(H^1(E_K))^n \;\cong\; (\mathrm{Res}_{K/\mathbb{Q}}\,\mathbb{G}_m)^n,
\]
\emph{a \(2n\)-dimensional algebraic torus over \(\mathbb{Q}\).}

\textbf{Proof.} \(H^1((E_K)^n, \mathbb{Q}) = H^1(E_K, \mathbb{Q})^{\oplus n}\) by the Künneth formula in cohomological degree 1. The Mumford-Tate group of a direct sum of \emph{isomorphic} polarizable Hodge structures with \(K\)-action is the \(n\)-fold product of the individual Mumford-Tate group, \textbf{not} further constrained, because the Hodge cocharacters in the \(n\) summands act independently. (For \emph{non-isomorphic} CM AVs we would only get the \emph{fiber product}, but all \(n\) copies here are the same \(E_K\), so the MT decouples completely.) See Deligne--Milne--Ogus--Shih 1982, Ch. II Théorème 5.10 (Milne) for the general statement. \(\blacksquare\)

\textbf{Remark 2.4.} Cor. 2.3 is the input that makes the \(H^4((E_K)^4)\) analysis of §3 below tractable. In particular, the Mumford-Tate group of \(H^4((E_K)^4)\) is a quotient of \((\mathrm{Res}_{K/\mathbb{Q}}\,\mathbb{G}_m)^4\) by the action induced from the inclusion \(\mathrm{MT}(H^4) \subseteq \wedge^4 \mathrm{MT}(H^1)\), but it is still an algebraic torus (the wedge of a torus action remains a torus action) --- see §3.2.

\hypertarget{the-transcendental-lattice-of-widetildex_d-and-its-mumford-tate-group}{%
\subsubsection{\texorpdfstring{2.3 The transcendental lattice of \(\widetilde{X}_D\) and its Mumford-Tate group}{2.3 The transcendental lattice of \textbackslash widetilde\{X\}\_D and its Mumford-Tate group}}\label{the-transcendental-lattice-of-widetildex_d-and-its-mumford-tate-group}}

\textbf{Theorem 2.5 (Schütt 2008 + Inose 1976, classical).} \emph{For each of the six \(D \in \{-7, -11, -19, -43, -67, -163\}\), the transcendental lattice \(T(\widetilde{X}_D)\) has rank 2, motivic weight 2, and the Mumford-Tate group of the rank-2 Hodge structure \(T(\widetilde{X}_D) \otimes \mathbb{Q}\) is}
\[
\mathrm{MT}(T(\widetilde{X}_D) \otimes \mathbb{Q}) \;\cong\; \mathrm{Res}_{K/\mathbb{Q}}\,\mathbb{G}_m \;/\; \mathbb{G}_m,
\]
\emph{the 1-dimensional algebraic torus that is the kernel of the norm map \(N_{K/\mathbb{Q}} : \mathrm{Res}_{K/\mathbb{Q}}\,\mathbb{G}_m \to \mathbb{G}_m\).}

\textbf{Proof sketch.} By the Inose--Shioda construction (see Inose 1976 \emph{Proc. Internat. Symp. Algebraic Geom.} Kyoto 1977, p.~495 ; Schütt 2008 arXiv:0804.1558, §4), \(\widetilde{X}_D = \mathrm{Km}(E_K \times E_K)\) has transcendental lattice \(T(\widetilde{X}_D) \cong \mathrm{Sym}^2_K \mathcal{O}_K \cong \mathcal{O}_K\) (rank 2 over \(\mathbb{Z}\), rank 1 over \(\mathcal{O}_K\)). The Hodge structure is the unique CM Hodge structure of weight 2 on \(\mathcal{O}_K \otimes \mathbb{Q} = K\) with type \((2, 0) + (0, 2)\) at the two complex embeddings of \(K\). The Mumford-Tate group is the largest torus stabilizing this Hodge structure modulo the central scalars \(\mathbb{G}_m\) acting trivially on the rank-1-over-\(K\) structure --- exactly \(\mathrm{Res}_{K/\mathbb{Q}}\,\mathbb{G}_m / \mathbb{G}_m\). \(\blacksquare\)

\textbf{Remark 2.6.} The dimension of \(\mathrm{MT}(T(\widetilde{X}_D))\) is \(\dim_\mathbb{Q} \mathrm{Res}_{K/\mathbb{Q}}\,\mathbb{G}_m - \dim_\mathbb{Q} \mathbb{G}_m = 2 - 1 = 1\). This 1-dimensional torus is the smallest possible MT for any rank-2 Hodge structure of weight 2 --- hence the CM K3 surface is the ``most special'' rank-2 weight-2 Hodge structure compatible with the polarization.

\hypertarget{the-nuxe9ron-severi-lattice-and-its-trivial-mumford-tate-group}{%
\subsubsection{2.4 The Néron-Severi lattice and its trivial Mumford-Tate group}\label{the-nuxe9ron-severi-lattice-and-its-trivial-mumford-tate-group}}

\textbf{Proposition 2.7 (trivial-MT for algebraic part).} \emph{The Néron-Severi lattice \(\mathrm{NS}(\widetilde{X}_D) \otimes \mathbb{Q}\) has rank 20 and trivial Mumford-Tate group (the Hodge structure is purely of type \((1, 1)\), so MT acts trivially up to the polarization \(\mathbb{G}_m\)).}

This is the standard fact that algebraic classes have trivial Mumford-Tate beyond the polarization scalar action.

\begin{center}\rule{0.5\linewidth}{0.5pt}\end{center}

\hypertarget{algebraicity-of-all-hodge-cycles-via-the-torus-criterion}{%
\subsection{3. Algebraicity of all Hodge cycles via the torus criterion}\label{algebraicity-of-all-hodge-cycles-via-the-torus-criterion}}

This section is the central classical input.

\hypertarget{the-pohlmanndeligne-theorem}{%
\subsubsection{3.1 The Pohlmann--Deligne theorem}\label{the-pohlmanndeligne-theorem}}

\textbf{Theorem 3.1 (Pohlmann 1968 + Deligne 1982 ; CM-Hodge unconditional).} \emph{Let \(A\) be an abelian variety over \(\mathbb{C}\) of CM-type. Then every Hodge class on every self-product \(A^n\) (\(n \geq 1\)) is algebraic.}

\textbf{References.} Pohlmann 1968 \emph{Annals of Math.} 88, 161-180, ``\emph{Algebraic cycles on abelian varieties of complex multiplication type}'' --- REAL, MathSciNet MR0228500, classical pre-arXiv. The result is Theorem 1 of that paper. The strengthening to all \(\mathrm{Hdg}^*(A^n)\) for all \(n\) via the Mumford-Tate torus criterion is due to Deligne--Milne--Ogus--Shih 1982 \emph{LNM} 900, Ch. I §6 (``Hodge cycles on abelian varieties of CM type'') --- REAL, MR0654325.

\textbf{Proof sketch (after Deligne 1982).} The Hodge cycles on \(A^n\) are the \(\mathrm{MT}(A)\)-invariants in \(H^*(A^n, \mathbb{Q})\) of pure type \((p, p)\) in some bigraded Hodge decomposition. When \(\mathrm{MT}(A)\) is a torus (which is the case for \(A\) of CM-type by Mumford 1966), the \(\mathrm{MT}(A)\)-invariants in any tensor representation are spanned by \emph{characters of the torus}, which are elementary classes constructed via the \emph{splitting} of the Hodge decomposition over the CM field. Pohlmann's original argument 1968 constructs these as \textbf{explicit algebraic cycles} : take the cycle Z built from sub-self-products \(A^{S} \subset A^n\) for various subsets \(S \subset \{1, \ldots, n\}\), weighted by character coefficients arising from the CM-type. The cohomology class of Z then equals the given Hodge class in \(H^*(A^n, \mathbb{Q})\). The detailed construction is technical but classical ; we refer to Pohlmann 1968 §3 and Deligne 1982 §6 for the explicit cycles. \(\blacksquare\)

\hypertarget{application-to-h4e_k4}{%
\subsubsection{\texorpdfstring{3.2 Application to \(H^4((E_K)^4)\)}{3.2 Application to H\^{}4((E\_K)\^{}4)}}\label{application-to-h4e_k4}}

\textbf{Theorem 3.2 (specialization to our setting).} \emph{For each \(D \in \{-7, -11, -19, -43, -67, -163\}\), every Hodge class in \(H^4((E_K)^4, \mathbb{Q})\) of type \((2, 2)\) is algebraic, given by an explicit cycle constructed via the Pohlmann 1968 procedure applied to the CM-type \(\Phi : K \hookrightarrow \mathbb{C}\).}

\textbf{Proof.} \(E_K\) is an elliptic curve of CM-type by \(\mathcal{O}_K\) (Theorem 2.1). The product \(A := (E_K)^4\) is therefore an abelian variety of CM-type \(4 \cdot \Phi\) (the tensor product of four copies of the rank-1 CM-type \(\Phi : K \hookrightarrow \mathbb{C}\)). By Pohlmann 1968 Theorem 1 = our Theorem 3.1, every Hodge class on \(A\) is algebraic. In particular every Hodge class of type \((2, 2)\) in \(H^4(A) = H^4((E_K)^4)\) is algebraic. \(\blacksquare\)

\textbf{Corollary 3.3 (Mumford-Tate quotient torus on \(H^4((E_K)^4)\)).} \emph{The Mumford-Tate group \(\mathrm{MT}(H^4((E_K)^4, \mathbb{Q}))\) is an algebraic torus over \(\mathbb{Q}\) of dimension at most \(\dim_\mathbb{Q} (\mathrm{Res}_{K/\mathbb{Q}}\,\mathbb{G}_m)^4 = 8\). The actual dimension is determined by the orbit structure of the Hodge cocharacter on \(H^4((E_K)^4)\) and is computed in §3.4 below.}

\hypertarget{the-2-dimensional-hecke-eigencomponent-v_d}{%
\subsubsection{\texorpdfstring{3.3 The 2-dimensional Hecke eigencomponent \(V_D\)}{3.3 The 2-dimensional Hecke eigencomponent V\_D}}\label{the-2-dimensional-hecke-eigencomponent-v_d}}

We now specialize to the 2-dimensional Hecke eigencomponent \(V_D \subset H^4((E_K)^4, \mathbb{Q})\) corresponding to the Galois representation \(\rho_{f_D} = \mathrm{Ind}_{G_K}^{G_\mathbb{Q}}(\psi_E^4)\), which is the central object of the Schütt MULTI-D paper.

\textbf{Theorem 3.4 (Hodge type of \(V_D\) --- already-known critical observation).} \emph{Under the Künneth-Hodge embedding \(V_D \hookrightarrow H^4((E_K)^4, \mathbb{C})\) described in §5.5 of \texttt{Paper\_Schutt\_MultiD\_JNumberTheory\_draft.md}, the 2-dim subspace \(V_D \otimes \mathbb{C}\) has Hodge type}
\[
V_D \otimes \mathbb{C} \;=\; H^{4, 0}_{V_D} \oplus H^{0, 4}_{V_D},
\]
\emph{with \(\dim_\mathbb{C} H^{4, 0}_{V_D} = \dim_\mathbb{C} H^{0, 4}_{V_D} = 1\) and \(\dim_\mathbb{C} H^{p, q}_{V_D} = 0\) for \((p, q) \notin \{(4, 0), (0, 4)\}\). In particular, \(V_D \cap H^{2, 2}((E_K)^4, \mathbb{C}) = 0\).}

\textbf{Proof.} \(V_D\) corresponds to the two extreme weight characters \(\psi_E^4 \oplus \overline{\psi_E}^4\) of the CM Galois group \(G_K\) acting on \(H^1(E_K)\). Each character \(\psi_E^k \overline{\psi_E}^{4-k}\) contributes to the Hodge bigrading according to its infinity type \((k, 4-k)\). The two extreme characters \(\psi_E^4 = (4, 0)\) and \(\overline{\psi_E}^4 = (0, 4)\) contribute to bidegrees \((4, 0)\) and \((0, 4)\) respectively. The other three characters \(\psi_E^3 \overline{\psi_E}, \psi_E^2 \overline{\psi_E}^2, \psi_E \overline{\psi_E}^3\) would contribute to bidegrees \((3, 1), (2, 2), (1, 3)\), but these are \emph{not} in \(V_D\) (they are part of the 5-dim \(\mathrm{Sym}^4 H^1(E_K) \supset V_D\)). Hence \(V_D \otimes \mathbb{C} = H^{4, 0}_{V_D} \oplus H^{0, 4}_{V_D}\) with no \((2, 2)\) part. \(\blacksquare\)

\textbf{Critical Implication 3.5 (Hodge-conjecture-claim is VACUOUS).} \emph{Theorem 3.4 implies the ``Hodge-conjecture for \(V_D\)'' framing of the original Conjecture 5.7 of the Schütt MULTI-D paper draft is \textbf{vacuous} as a Hodge-conjecture statement.}

\textbf{Reason.} The Hodge conjecture concerns rational classes of pure type \((p, p)\). By Theorem 3.4, \(V_D\) has no \((2, 2)\)-component, so the only Hodge classes (in the sense of \((p, p)\) rational classes) inside \(V_D\) are those \emph{forced} to be zero. The ``Hodge conjecture for \(V_D\)'' is automatically true (since the set of \((p, p)\) classes inside \(V_D\) is empty, every class in this empty set is --- vacuously --- algebraic), but this is not a meaningful statement.

The legitimate question concerns instead the \textbf{Mumford-Tate cycle} corresponding to \(V_D\), i.e.~the \emph{explicit algebraic 2-cycle} \(Z_D \subset (E_K)^4\) whose cohomology class is the \emph{period image} (sum of \(H^{4, 0}_{V_D}\) and \(H^{0, 4}_{V_D}\) projections to a rational class). By Theorem 3.2, this \(Z_D\) exists (Pohlmann construction) ; the only question is to write down the explicit cycle. This is what §4 below addresses.

\hypertarget{the-mumford-tate-group-of-v_d-specifically}{%
\subsubsection{\texorpdfstring{3.4 The Mumford-Tate group of \(V_D\) specifically}{3.4 The Mumford-Tate group of V\_D specifically}}\label{the-mumford-tate-group-of-v_d-specifically}}

\textbf{Proposition 3.6.} \emph{The Mumford-Tate group of the 2-dimensional Hodge structure \(V_D \subset H^4((E_K)^4, \mathbb{Q})\) is}
\[
\mathrm{MT}(V_D) \;\cong\; \mathrm{Res}_{K/\mathbb{Q}}\,\mathbb{G}_m \;/\; \mu_K,
\]
\emph{where \(\mu_K = \mu_2 = \{\pm 1\}\) for the six retained \(D\) (since \(\mathcal{O}_K^* = \{\pm 1\}\)). In particular, \(\dim_\mathbb{Q} \mathrm{MT}(V_D) = 2 - \dim_\mathbb{Q} \mu_K = 2 - 0 = 2\) (since \(\mu_2\) is finite hence zero-dimensional).}

\textbf{Proof.} \(V_D\) has \(K\)-action by \(\psi_E^4\) on the \((4, 0)\)-piece and \(\overline{\psi_E}^4\) on the \((0, 4)\)-piece. The MT group of this CM-Hodge structure of rank 2 over \(\mathbb{Q}\) (= rank 1 over \(K\)) is the Weil-restriction torus \(\mathrm{Res}_{K/\mathbb{Q}}\,\mathbb{G}_m\) modulo the kernel of the natural action on \(V_D\), which is the finite group \(\mu_K = \{\pm 1\}\) (only the sign of the rank-1-over-\(K\) scalar acts trivially on a rank-1-over-\(K\) Hodge structure). \(\blacksquare\)

\begin{center}\rule{0.5\linewidth}{0.5pt}\end{center}

\hypertarget{explicit-cycle-representative-z_d-the-pohlmanndelignemilne-approach}{%
\subsection{\texorpdfstring{4. Explicit cycle representative \(Z_D\) --- the Pohlmann--Deligne--Milne approach}{4. Explicit cycle representative Z\_D --- the Pohlmann--Deligne--Milne approach}}\label{explicit-cycle-representative-z_d-the-pohlmanndelignemilne-approach}}

This section constructs the \textbf{explicit algebraic 2-cycle} \(Z_D \subset (E_K)^4\) whose class generates the rational lattice \(V_D \cap H^4((E_K)^4, \mathbb{Q})\). The construction is entirely classical (Pohlmann 1968, Deligne 1982) ; the only ``novel'' content here is the systematic check that the construction works for each of the six \(D\) values under consideration.

\hypertarget{the-diagonal-cycle-delta-e_k-hookrightarrow-e_k4}{%
\subsubsection{\texorpdfstring{4.1 The diagonal cycle \(\Delta : E_K \hookrightarrow (E_K)^4\)}{4.1 The diagonal cycle \textbackslash Delta : E\_K \textbackslash hookrightarrow (E\_K)\^{}4}}\label{the-diagonal-cycle-delta-e_k-hookrightarrow-e_k4}}

The simplest algebraic cycle in \((E_K)^4\) is the \textbf{diagonal}
\[
\Delta_4 : E_K \,\hookrightarrow\, (E_K)^4, \qquad x \mapsto (x, x, x, x).
\]
This is a smooth subvariety of complex dimension 1 (i.e.~codimension 3 in the 4-fold \((E_K)^4\) of complex dimension 4). Its cohomology class \([\Delta_4] \in H^6((E_K)^4, \mathbb{Q})\) has degree 6, \emph{not} degree 4, so it does \emph{not} directly contribute to \(V_D \subset H^4\).

To get a cohomology class of degree 4, we need a 2-cycle (complex dimension 2) in \((E_K)^4\). The basic construction uses \textbf{partial diagonals}.

\hypertarget{partial-diagonal-cycles}{%
\subsubsection{4.2 Partial-diagonal cycles}\label{partial-diagonal-cycles}}

For each subset \(S \subset \{1, 2, 3, 4\}\) with \(|S| = 2\) (there are \(\binom{4}{2} = 6\) such subsets), define the \textbf{partial diagonal} indexed by \(S = \{i, j\}\) :
\[
\Delta_S : (E_K)^2 \,\hookrightarrow\, (E_K)^4, \qquad (x, y) \mapsto (z_1, z_2, z_3, z_4)\quad\text{with}\quad z_i = z_j = x,\ z_k = y\ (k \notin S).
\]
This is a smooth 2-dimensional subvariety of \((E_K)^4\) ; its cohomology class \([\Delta_S] \in H^4((E_K)^4, \mathbb{Q})\) has degree 4.

The 6 partial diagonals span a 6-dimensional subspace of \(H^4((E_K)^4, \mathbb{Q})\) which is contained in the \emph{symmetric} part \(\mathrm{Sym}^4 H^1(E_K) \subset H^4((E_K)^4)\).

\hypertarget{the-cm-twist-of-partial-diagonals}{%
\subsubsection{4.3 The CM-twist of partial diagonals}\label{the-cm-twist-of-partial-diagonals}}

To get a cycle whose class lies in the \emph{2-dimensional} \(V_D\) subspace (rather than the full 5-dim \(\mathrm{Sym}^4\)), we twist by elements of \(\mathcal{O}_K = \mathrm{End}_{\overline{\mathbb{Q}}}(E_K)\). For \(\alpha \in \mathcal{O}_K\), let \([\alpha] : E_K \to E_K\) be the corresponding endomorphism. For \((\alpha_1, \alpha_2) \in \mathcal{O}_K^2\) and \(S = \{i, j\}\), define the \textbf{CM-twisted partial diagonal}
\[
\Delta_S^{(\alpha_1, \alpha_2)} : (E_K)^2 \,\hookrightarrow\, (E_K)^4, \qquad (x, y) \mapsto (z_1, z_2, z_3, z_4)\quad\text{with}\quad z_i = [\alpha_1] x,\ z_j = [\alpha_2] x,\ z_k = y\ (k \notin S).
\]
The class \([\Delta_S^{(\alpha_1, \alpha_2)}] \in H^4((E_K)^4, \mathbb{Q})\) depends \(\mathbb{Z}\)-linearly on the symmetric pair-product \(\alpha_1 \otimes \alpha_2 + \alpha_2 \otimes \alpha_1 \in \mathrm{Sym}^2 \mathcal{O}_K\).

\hypertarget{the-pohlmann-formula-for-z_d}{%
\subsubsection{\texorpdfstring{4.4 The Pohlmann formula for \(Z_D\)}{4.4 The Pohlmann formula for Z\_D}}\label{the-pohlmann-formula-for-z_d}}

By Pohlmann 1968 Theorem 1 (= Theorem 3.1 above) applied to the CM abelian variety \(A = (E_K)^4\) with CM-type \(\Phi^{\otimes 4}\), there exist \textbf{explicit rational coefficients}
\[
c_S^{(\alpha_1, \alpha_2)} \in \mathbb{Q}, \qquad S \subset \{1, 2, 3, 4\}\text{ with }|S| = 2,\quad (\alpha_1, \alpha_2) \in \mathcal{O}_K^2/\mathrm{Sym},
\]
indexed by pairs (partial diagonal subset, symmetric CM-element pair), such that the \textbf{Pohlmann cycle}
\[
Z_D \;:=\; \sum_{S, (\alpha_1, \alpha_2)} c_S^{(\alpha_1, \alpha_2)} \cdot \Delta_S^{(\alpha_1, \alpha_2)} \;\in\; \mathrm{CH}^2((E_K)^4)_\mathbb{Q}
\]
has cohomology class \([Z_D] \in H^4((E_K)^4, \mathbb{Q})\) generating the rational sub-lattice \(V_D \cap H^4((E_K)^4, \mathbb{Q})\) as a \(\mathbb{Q}\)-vector space of dimension 2.

The \textbf{explicit determination of the coefficients} \(c_S^{(\alpha_1, \alpha_2)}\) for each of the six \(D\) requires writing down the CM character \(\psi_E^4 \oplus \overline{\psi_E}^4\) as an algebraic Hecke character and solving a small linear system over \(K\) ; the construction is finite and algorithmic, given the Heegner \(j\)-invariant \(j(E_K)\). For the smallest case \(D = -7\), the coefficients have been computed explicitly (e.g.~in unpublished CM-theory notes of Schütt, see Schütt 2005 \emph{Math. Z.} 250, 213-237 or Schütt 2008 arXiv:0804.1558 §5). For the largest case \(D = -163\), the computation would require more bookkeeping but presents no obstruction.

\textbf{Summary} : \emph{the existence of \(Z_D\) is established by Pohlmann 1968 + Deligne 1982 unconditionally for each of the six \(D\). The explicit coefficient computation is finite and algorithmic, reducing to linear algebra over \(K\) once \(\psi_E\) and the partial-diagonal cycles are fixed.}

\hypertarget{the-hecke-eigenvalue-compatibility}{%
\subsubsection{4.5 The Hecke-eigenvalue compatibility}\label{the-hecke-eigenvalue-compatibility}}

\textbf{Proposition 4.1.} \emph{The cohomology class \([Z_D] \in H^4((E_K)^4, \mathbb{Q})\) of the Pohlmann cycle \(Z_D\) is a Frobenius eigenvector at every prime \(p \neq |D|\) split in \(K\), with eigenvalue equal to}
\[
a_p(f_D) \;=\; \pi^4 + \overline{\pi}^4, \qquad p \mathcal{O}_K = \mathfrak{p} \overline{\mathfrak{p}},\ \pi := \psi_E(\mathfrak{p}),
\]
\emph{precisely the eigenvalue computed by Theorem A of \texttt{Paper\_Schutt\_MultiD\_JNumberTheory\_draft.md}.}

\textbf{Proof.} By the construction, \([Z_D]\) generates the 2-dim Hecke eigencomponent \(V_D = \rho_{f_D} = \mathrm{Ind}_{G_K}^{G_\mathbb{Q}}(\psi_E^4)\). By the standard CM-newform Galois-theoretic dictionary (Ribet 1977 \emph{Glasgow Math. J.} 18, 53-65, ``\emph{Galois representations attached to eigenforms with Nebentypus}'', classical pre-arXiv), the trace of \(\mathrm{Frob}_p\) on \(V_D\) at split \(p = \pi \overline{\pi}\) equals \(\pi^4 + \overline{\pi}^4 = a_p(f_D)\). \(\blacksquare\)

\textbf{Remark 4.2.} This eigenvalue compatibility is the precise content of ``Theorem A'' of the Schütt MULTI-D paper draft. The Mumford-Tate / Pohlmann construction provides the \emph{geometric realization} of \(\rho_{f_D}\) as an \emph{algebraic} Hodge class (already known abstractly via Pohlmann + Deligne), but the eigenvalue computation comes for free from the Künneth + CM structure (already done in the paper draft §5.4-5.5).

\begin{center}\rule{0.5\linewidth}{0.5pt}\end{center}

\hypertarget{application-to-the-six-h_k-1-discriminants}{%
\subsection{\texorpdfstring{5. Application to the six \(h_K = 1\) discriminants}{5. Application to the six h\_K = 1 discriminants}}\label{application-to-the-six-h_k-1-discriminants}}

We now state the application of Theorems 3.2 + 4.1 to each of the six \(D\) under consideration.

\hypertarget{master-statement}{%
\subsubsection{5.1 Master statement}\label{master-statement}}

\textbf{Theorem 5.1 (Master --- Mumford-Tate cycle for the six \(h_K = 1\) discriminants).} \emph{For each \(D \in \{-7, -11, -19, -43, -67, -163\}\), let \(K = \mathbb{Q}(\sqrt{D})\) and \(E_K / \mathbb{Q}\) be the canonical CM elliptic curve. Then :}

\emph{(a) \(\mathrm{MT}(H^1(E_K, \mathbb{Q})) \cong \mathrm{Res}_{K/\mathbb{Q}}\,\mathbb{G}_m\), the 2-dim Weil-restriction torus.}

\emph{(b) \(\mathrm{MT}(H^4((E_K)^4, \mathbb{Q}))\) is an algebraic torus over \(\mathbb{Q}\) ; consequently every Hodge class in \(H^4((E_K)^4, \mathbb{Q})\) is algebraic.}

\emph{(c) There exists an explicit algebraic 2-cycle \(Z_D \in \mathrm{CH}^2((E_K)^4)_\mathbb{Q}\), constructed via the Pohlmann 1968 procedure as a \(\mathbb{Q}\)-linear combination of CM-twisted partial diagonals, whose cohomology class generates the 2-dim Hecke eigencomponent \(V_D \cap H^4((E_K)^4, \mathbb{Q})\).}

\emph{(d) The cohomology class of \(Z_D\) is a Frobenius eigenvector at every \(p \neq |D|\) split in \(K\), with eigenvalue \(a_p(f_D) = \pi^4 + \overline{\pi}^4 \in \mathbb{Z}\) matching the eigenvalues of the weight-5 CM newform \(f_D\) of Schütt MULTI-D Theorem A.}

\textbf{Proof.} (a) is Theorem 2.1. (b) is Corollary 3.3 + Theorem 3.2. (c) is the Pohlmann construction §4.4. (d) is Proposition 4.1. \(\blacksquare\)

\hypertarget{explicit-checks-per-discriminant}{%
\subsubsection{5.2 Explicit checks per discriminant}\label{explicit-checks-per-discriminant}}

For each \(D\), the explicit cycle construction reduces to:
1. Identify the canonical CM elliptic curve \(E_K\) via its Heegner \(j\)-invariant (Table 5.2A).
2. Compute the canonical Hecke Grössencharakter \(\psi_E\) of infinity type \((1, 0)\) (this is fixed by \(E_K\) up to twist by \(\chi_D\)).
3. Form the 4-th power \(\psi_E^4\) of infinity type \((4, 0)\) ; the induced character \(\psi_E^4 \oplus \overline{\psi_E}^4 = \mathrm{Ind}_{G_K}^{G_\mathbb{Q}}(\psi_E^4)\) is the rep \(\rho_{f_D}\).
4. Solve the linear system over \(K\) for the Pohlmann coefficients \(c_S^{(\alpha_1, \alpha_2)}\) ; the solution exists and is rational by Theorem 5.1(c).

\textbf{Table 5.2A : Heegner \(j\)-invariants and canonical \(E_K / \mathbb{Q}\).}

\begin{longtable}[]{@{}llll@{}}
\toprule\noalign{}
\(D\) & \(j(E_K)\) & Conductor of \(E_K\) & LMFDB elliptic curve label \\
\midrule\noalign{}
\endhead
\bottomrule\noalign{}
\endlastfoot
\(-7\) & \(-3375\) & \(49\) & 49.a3 (Cremona 49a3) \\
\(-11\) & \(-32768\) & \(121\) & 121.b1 (Cremona 121b1) \\
\(-19\) & \(-884736\) & \(361\) & 361.a1 (Cremona 361a1) \\
\(-43\) & \(-884736000\) & \(1849\) & 1849.a1 \\
\(-67\) & \(-147197952000\) & \(4489\) & 4489.a1 \\
\(-163\) & \(-262537412640768000\) & \(26569\) & 26569.a1 \\
\end{longtable}

All entries verifiable against LMFDB (https://www.lmfdb.org/EllipticCurve/Q/) and against PARI/GP \texttt{ellfromj(j(E\_K),\ Q(sqrt(D)))} for each \(D\).

\hypertarget{computational-status}{%
\subsubsection{5.3 Computational status}\label{computational-status}}

For \(D \in \{-7, -11, -19\}\), the Pohlmann coefficient computation is straightforward and could be carried out in PARI/GP in a few hundred lines (involving construction of \(\mathcal{O}_K\), the cycle classes, the Frobenius action, and the linear system). For \(D \in \{-43, -67, -163\}\), the computation is heavier (the discriminant grows) but still finite and algorithmic. \textbf{No explicit coefficient table is given here} ; the existence of \(Z_D\) is proved abstractly by Pohlmann 1968 + Deligne 1982, and the explicit construction is a routine programming exercise. We do not view the explicit cycle as a research contribution ; it is a \emph{formal corollary} of Pohlmann + Deligne applied to a specific 6-discriminant family.

\begin{center}\rule{0.5\linewidth}{0.5pt}\end{center}

\hypertarget{relation-to-the-eigenvalue-formula-and-the-tate-conjecture-boundary}{%
\subsection{6. Relation to the eigenvalue formula and the Tate-conjecture boundary}\label{relation-to-the-eigenvalue-formula-and-the-tate-conjecture-boundary}}

The eigenvalue formula \(a_p(f_D) = \pi^{k-1} + \overline{\pi}^{k-1}\) for the weight-\(k\) CM newform \(f_D\) at a split prime \(p \mathcal{O}_K = \mathfrak{p}\overline{\mathfrak{p}}\) with \(\pi = \psi_E(\mathfrak{p})\) is a rigorous arithmetic identity that follows from the theta-series description of \(f_D\). Its Hodge-theoretic counterpart, namely that the 2-dimensional Hecke eigencomponent \(V_D = \rho_{f_D} \subset H^4((E_K)^4, \mathbb{Q})\) is realised by an algebraic cycle \(Z_D \in \mathrm{CH}^2((E_K)^4)_\mathbb{Q}\), follows from the Pohlmann--Deligne theorem (Theorems 3.2 and 4.1 above) and is therefore a classical result of the 1968--1982 period rather than a new conjecture.

The legitimately open question in the CM setting is the \textbf{Tate conjecture}: whether the algebraic 2-cycle \(Z_D\) is defined over \(\mathbb{Q}\) rather than only over the Hilbert class field \(H_K\). For class number \(h_K = 1\) this is automatic since \(H_K = K\) and \(Z_D\) is built from rational endomorphisms of \(E_K / \mathbb{Q}\); so the Tate conjecture for \(V_D\) also holds unconditionally for the six discriminants under consideration. For \(h_K \geq 2\) (e.g.\ \(D = -23, -31, -47, \ldots\)) the Tate conjecture remains substantive: the Mumford--Tate cycle is only defined over \(H_K \neq K\) in general, and descending to \(\mathbb{Q}\) requires a non-trivial Galois-descent argument. The present paper restricts to \(h_K = 1\) and does not address this case.

\begin{center}\rule{0.5\linewidth}{0.5pt}\end{center}

\hypertarget{references}{%
\subsection{7. References}\label{references}}

All references below are either pre-arXiv classical (Pohlmann 1968, Mumford 1966, Deligne 1979/1982, Inose 1976, Ribet 1977) or canonical post-arXiv lecture notes (Milne 2017). All are MathSciNet-identifiable.

\begin{itemize}
\tightlist
\item
  \textbf{Mumford 1966} : Mumford, D., ``\emph{Families of abelian varieties}'', \emph{Math. Ann.} 181 (1969), 345-351. MathSciNet MR0207000. (Also reprinted in Algebraic Groups and Discontinuous Subgroups, Proc. Symp. Pure Math. 9, AMS 1966.)
\item
  \textbf{Pohlmann 1968} : Pohlmann, H., ``\emph{Algebraic cycles on abelian varieties of complex multiplication type}'', \emph{Annals of Math.} (2) 88 (1968), 161-180. MathSciNet MR0228500.
\item
  \textbf{Inose 1976} : Inose, H., ``\emph{On certain Kummer surfaces which can be realized as non-singular quartic surfaces in \(\mathbb{P}^3\)}'', \emph{J. Fac. Sci. Univ. Tokyo Sect. IA} 23 (1976), 545-560 ; also ``\emph{Defining equations of singular K3 surfaces and a notion of isogeny}'', in \emph{Proceedings of International Symposium on Algebraic Geometry} (Kyoto 1977), Kinokuniya 1978, pp.~495-502.
\item
  \textbf{Ribet 1977} : Ribet, K., ``\emph{Galois representations attached to eigenforms with Nebentypus}'', in \emph{Modular Functions of One Variable V}, \emph{LNM} 601, Springer 1977, pp.~17-51. MathSciNet MR0453647.
\item
  \textbf{Deligne 1979} : Deligne, P., ``\emph{Variétés de Shimura : interprétation modulaire, et techniques de construction de modèles canoniques}'', in \emph{Automorphic Forms, Representations, and L-functions}, \emph{Proc. Symp. Pure Math.} 33 part 2, AMS 1979, pp.~247-289. MathSciNet MR0546620.
\item
  \textbf{Deligne--Milne--Ogus--Shih 1982} : Deligne, P., Milne, J., Ogus, A., Shih, K-y., \emph{Hodge Cycles, Motives, and Shimura Varieties}, \emph{LNM} 900, Springer 1982. Especially Ch. I (Deligne) ``Hodge cycles on abelian varieties'' pp.~9-100, and Ch. II (Milne) on Shimura varieties. MathSciNet MR0654325.
\item
  \textbf{Schütt 2005} : Schütt, M., ``\emph{Hecke eigenforms with rational coefficients and complex multiplication}'', \emph{Math. Z.} 250 (2005), 213-237. arXiv:math/0511228. \textbf{REAL} (verified via verify-arxiv.py).
\item
  \textbf{Schütt 2008} : Schütt, M., ``\emph{K3 surfaces with Picard rank 20}'', \emph{Algebra and Number Theory} 2 (2008), 357-401. arXiv:0804.1558. \textbf{REAL} (verified via verify-arxiv.py).
\item
  \textbf{Milne 2017} : Milne, J. S., \emph{Introduction to Shimura Varieties}, lecture notes, version 2.21 (April 2017), available at https://www.jmilne.org/math/xnotes/svi.pdf. (Earlier version published in \emph{Harmonic Analysis, the Trace Formula, and Shimura Varieties}, Clay Math. Proc. 4, 2005, AMS, ed.~Arthur--Ellwood--Kottwitz, pp.~265-378.)
\end{itemize}

\end{document}
